\par\SourceLeaf{435}\Needspace{7\baselineskip}\bigskip
\begin{center}
\pdfbookmark[2]{Gauss an Olbers, Göttingen, 2. December 1817}{pallas-brief-18171202}
\textsc{Gauss} an \textsc{Olbers}, Göttingen, 2. December 1817.
\end{center}
\noindent\hspace*{\parindent}. . . . .\quad Ich glaube, Ihnen schon einigemale von der Tafel für die Erleichterung der Berechnung des jedesmaligen Betrags der Pallasstörungen geschrieben zu haben. Seitdem ich die Beihülfe des Hrn. \textsc{Westphal} dabei entbehrte, habe ich an dem noch fehlenden Theil derselben nur langsam arbeiten mögen, seit einiger Zeit ist sie jedoch vollendet. Ich habe es nun für nöthig gehalten, die sämmtlichen bisher beobachteten 12 Gegenscheine von neuem mit Hülfe dieser Tafel mit Normalelementen zu vergleichen, die daraus entspringenden 23 Bedingungsgleichungen (in Einem Gegenschein taugt die Breite nichts, da bloss \textsc{David} sehr schlechte Declinationen beobachtete, No. V. von 1808) zu entwickeln und die 7 Correctionen (die der {\AstroSymbols ♃}-Masse eingeschlossen) daraus zu bestimmen. Diese mühsame Arbeit habe ich jetzt vorläufig (mit Vorbehalt einer kleinen Nachfeilung) vollendet. Die Übereinstimmung ist in der That zu bewundern, um so mehr, da ich die Saturnsstörungen nicht zugezogen habe, die Hr. \textsc{Nicolai} berechnet hat, und die die Übereinstimmung merklich \textls[180]{schlechter} machen würden. Die Correction der Jupitersmasse \((+\frac{1}{46})\) bestätigt sich nun immer mehr, die Summe der Quadrate der 12 Längenunterschiede, welche ohne diese Verbesserung nicht unter \(4064''\) heruntergebracht werden könnte, wird mit derselben bis auf \(241''\) gebracht, und nach meiner Wahrscheinlichkeitstheorie ist jetzt die Bestimmung der Jupitersmasse aus den Pallasstörungen etwa dreimal so genau, als aus den Saturnsbeobachtungen, aus denen bekanntlich \textsc{Laplace} ein mit dem meinigen in Widerspruch stehendes Resultat abgeleitet hat. Die Zuziehung der \textsc{Nicolai}schen Störungen vom Saturn würde übrigens die Correction der Jupitersmasse nicht abändern. Sehen Sie hier die übriggebliebenen Differenzen:
\begin{center}
\begingroup\small\setlength{\tabcolsep}{5pt}
\begin{tabular}{@{}lr r r r r@{}}
&&\multicolumn{1}{c}{Heliocentr.}&\multicolumn{1}{c}{Geocentr.}&\multicolumn{2}{c}{Differenz nach \textsc{Carlinis}}\\
&&\multicolumn{1}{c}{Länge}&\multicolumn{1}{c}{Breite}&\multicolumn{2}{c}{Bestimmung der Oppositionen}\\[3pt]
I.&1803&\(+0''{,}8\)&\(-1''{,}2\)&\(-5''{,}2\)&\(-30''{,}2\)\\
II.&1804&\(+4{,}0\)&\(-11{,}8\)&\(+4{,}5\)&\(-9{,}8\)\\
III.&1805&\(-2{,}1\)&\(-3{,}8\)&\(-3{,}7\)&\(+3{,}9\)\\
IV.&1807&\(-4{,}6\)&\(-2{,}3\)&\(+1{,}1\)&\(-4{,}5\)\\
V.&1808&\(-5{,}0\)&\(-16{,}5\mathpunct{::}\)&&
\end{tabular}
\endgroup
\end{center}
\SourceLeaf{436}
\begin{center}
\begingroup\small\setlength{\tabcolsep}{5pt}
\begin{tabular}{@{}lr r r r r@{}}
&&\multicolumn{1}{c}{Heliocentr.}&\multicolumn{1}{c}{Geocentr.}&\multicolumn{2}{c}{Differenz nach \textsc{Carlinis}}\\
&&\multicolumn{1}{c}{Länge}&\multicolumn{1}{c}{Breite}&\multicolumn{2}{c}{Bestimmung der Oppositionen}\\[3pt]
VI.&1809&\(+0''{,}7\)&\(-5''{,}1\)&&\\
VII.&1811&\(+2{,}4\)&\(-6{,}6\)&\(+2{,}9\)&\(+6{,}0\)\\
VIII.&1812&\(+10{,}7\)&\(+4{,}8\)&&\\
IX.&1813&\(-0{,}7\mathpunct{\raisebox{1ex}{.}}\)&\(-0{,}8\)&&\\
X.&1814&\(+1{,}8\)&\(-1{,}1\)&&\\
XI.&1816&\(-2{,}7\)&\(-8{,}5\)&&\\
XII.&1817&\(-6{,}5\)&\(+5{,}0\)&&
\end{tabular}
\endgroup
\end{center}
\par Diese Unterschiede sind um so unbedeutender, da die Bestimmung mehrerer Oppositionen von andern Astronomen, z. B. \textsc{Carlini}, öfters von der meinigen mehr abweicht, als die obigen Differenzen betragen, Da l’appétit vient en mangeant, so denke ich die Oppositionen II und VIII noch einmal aus den Beobachtungen mit Zuziehung meiner osculirenden Elemente berechnen zu lassen. Auch die Beobachtungen von 1802 denke ich noch discutiren zu lassen. Die Längen stimmen jetzt noch besser als die Breiten, was nicht zu verwundern ist, da jene heliocentrisch, diese geocentrisch sind.
\par
